% ============================================================
% 中文课程/实验报告模板（report-cn）
% ------------------------------------------------------------
% 基于 ctexart 文档类，默认使用 TeX Live 自带的 fandol 字体，
% 无需在系统中安装额外字体。编译引擎：xelatex。
% 修改指南：
%   1. 在"封面信息"处填写学校、课程、题目、姓名学号等；
%   2. 评分表如课程不要求，可整段删除（从"评分表"到 \clearpage）；
%   3. 正文按"目的/原理/步骤/结果/心得"组织，
%      不需要的章节直接删除；
%   4. 图、表、公式各有示例，替换文件名与数据即可。
% ============================================================
\documentclass[
  UTF8,
  fontset=fandol,        % TeX Live 自带开源字体，跨平台可编译
  zihao=-4,              % 正文字号：小四
  a4paper,
]{ctexart}

% ---- 常用宏包 ----
\usepackage{amsmath, amssymb}             % 数学公式
\usepackage{graphicx}                     % 插图
\usepackage{booktabs}                     % 三线表
\usepackage{array}                        % 表格列格式增强
\usepackage{geometry}
\geometry{left=2.5cm, right=2.5cm, top=2.6cm, bottom=2.6cm}
\usepackage{setspace}
\setstretch{1.4}                          % 正文 1.4 倍行距
\usepackage[
  colorlinks=true,
  linkcolor=black,
  citecolor=black,
  urlcolor=black,
]{hyperref}

% ---- 章节样式：传统中文报告风格"一、二、三、" ----
\ctexset{
  section = {
    format = \heiti\zihao{4},
    name = {,、},
    number = \chinese{section},
    aftername = {},
  },
  subsection = {
    format = \heiti\zihao{-4},
  },
}

% 图表标题样式
\usepackage{caption}
\captionsetup{font=small, labelfont=bf}

% ============================================================
% 封面信息：在这里填写报告信息
% ============================================================
\newcommand{\SchoolName}{示例大学}
\newcommand{\CourseName}{数据结构}
\newcommand{\ReportTitle}{实验二~二叉树及其应用}
\newcommand{\ReportAuthor}{张~三}
\newcommand{\StudentNumber}{2022010000}
\newcommand{\ClassName}{计算机 2201 班}
\newcommand{\Instructor}{李~四}
\newcommand{\ReportDate}{2026年5月20日}

\begin{document}

% ==================== 封面 ====================
\begin{titlepage}
  \begin{center}
    \vspace*{2em}
    {\zihao{2}\heiti \SchoolName}\\[1em]
    {\zihao{3}\heiti 课程实验报告}\\[5em]
    {\zihao{3}
      \renewcommand{\arraystretch}{1.8}
      \begin{tabular}{p{7em}p{16em}}
        {\heiti 课程名称：} & \CourseName \\
        {\heiti 实验题目：} & \ReportTitle \\
        {\heiti 姓~~~~名：} & \ReportAuthor \\
        {\heiti 学~~~~号：} & \StudentNumber \\
        {\heiti 班~~~~级：} & \ClassName \\
        {\heiti 指导教师：} & \Instructor \\
      \end{tabular}}\\[6em]
    {\zihao{3} \ReportDate}
  \end{center}
\end{titlepage}

% ==================== 评分表（课程不要求可整段删除） ====================
\begin{center}
  {\zihao{4}\heiti 评~分~表}
\end{center}
\vspace{0.5em}
\begin{center}
  \zihao{-4}
  \renewcommand{\arraystretch}{1.8}
  \begin{tabular}{|>{\centering\arraybackslash}p{8em}|p{22em}|>{\centering\arraybackslash}p{6em}|}
    \hline
    \textbf{评分项目} & \centering\textbf{评分要点} \tabularnewline
    \hline
    实验操作 & 规范完成实验步骤，能正确使用相关工具与环境 & \hspace{3em}/20 \\
    \hline
    结果正确性 & 实验结果正确、数据完整，达到预期目标 & \hspace{3em}/30 \\
    \hline
    报告质量 & 分析深入、条理清晰、图表规范、书写工整 & \hspace{3em}/40 \\
    \hline
    实验纪律 & 按时到场、独立完成、遵守机房规范 & \hspace{3em}/10 \\
    \hline
    \multicolumn{2}{|r|}{\textbf{总~分}} & \hspace{3em}/100 \\
    \hline
  \end{tabular}
\end{center}
\vspace{2em}
\begin{flushright}
  指导教师签字：\underline{\hspace{8em}} \quad
  批阅日期：\underline{\hspace{8em}}
\end{flushright}
\clearpage

% ==================== 正文 ====================

\section{实验目的}
% 逐条列出本实验要达到的目标
\begin{itemize}
  \item 理解二叉树的结构特性与存储实现；
  \item 掌握二叉树的遍历算法及其递归与非递归实现；
  \item 能够将二叉树应用于实际问题求解。
\end{itemize}

\section{实验原理}
% 简述实验涉及的理论知识，公式有编号示例
二叉树的先序遍历按照``根—左—右''的顺序访问结点。
设结点数为 $n$，树高为 $h$，则遍历的时间复杂度为 $O(n)$，
递归栈空间为 $O(h)$。

\section{实验内容与步骤}
% 分步描述实验做了什么，代码可粘贴关键片段并加注释
\begin{enumerate}
  \item 构造示例二叉树，采用二叉链表存储；
  \item 实现先序、中序、后序三种遍历；
  \item 对同一棵树分别输出遍历序列并核对结果；
  \item 统计不同规模随机树的遍历耗时。
\end{enumerate}

\section{实验结果与分析}
% 展示运行结果（表格或截图），并对结果进行解释
实验测得不同规模下的遍历耗时如表~\ref{tab:time} 所示，
结果与理论复杂度 $O(n)$ 一致。

\begin{table}[!htbp]
  \centering
  \caption{随机二叉树遍历耗时测量结果}
  \label{tab:time}
  \begin{tabular}{lcccc}
    \toprule
    结点数 $n$ & $10^3$ & $10^4$ & $10^5$ & $10^6$ \\
    \midrule
    耗时（ms） & 0.8 & 8.6 & 92.4 & 1043.1 \\
    \bottomrule
  \end{tabular}
\end{table}

% 插图示例：把截图文件放在本目录，替换文件名即可（支持 PNG/JPG/PDF）
% \begin{figure}[!htbp]
%   \centering
%   \includegraphics[width=0.7\linewidth]{result.png}
%   \caption{程序运行结果截图}
%   \label{fig:result}
% \end{figure}

\section{问题与心得}
% 记录调试中遇到的问题及解决办法，总结收获
实验中曾遇到递归深度过大导致栈溢出的问题，
改为显式栈的非递归实现后解决。通过本实验加深了
对树形结构与其遍历算法的理解。

\end{document}
